% Nouveau EGA V (cote 25): archive cote 25, site page 180 (PDF page 181).
% English translation of the checked French annotated master.
% AI-translated by GPT-6 Astra (OpenAI, ChatGPT Pro mode).
\scanpage{25/180}
\mspage{\hand{7}}

faithfully flat over S, which will prove (ii)$\hand{\Longrightarrow}$(iii bis) in 16.9.

Proof of 16.11. Since both members are smooth over S,
it suffices to prove flatness fiber by fiber, which reduces us to the
case where S is the spectrum of an alg closed field k. One can \var{\xout{then}}{}{typewriter strike-out}{ega25_25_180-v0001} obviously suppose
$\underline{E}_o=\underline{\unsure{B|E}}$, so $(**)$ takes the form
$(***)$   $Gl(\underline{E})\hand{\longrightarrow}\underline{Q}(\underline{E})^*$ ,
this morphism being obtained from the natural actions of the group scheme
$Gl(\underline{E})$ on $\underline{Q}(\underline{E})^*$ , by $u\hand{\rightsquigarrow} u(q_o)$, where $q_o$ is a section of $\underline{Q}(\underline{E})^*$ corresponding
to the standard form $Q_o$ . Now by the \unsure{gener}\unsure{ic} flatness theorem\hand{,}
since $\underline{Q}(\underline{E})^*$ is smooth over k, hence reduced, there is a dense open U of
$\underline{Q}(\underline{E})^*$ over which the \unsure{preceding} morphism is flat. For every \var{\margin{extract a lemma on \unsure{morphisms} of \unsure{h}omogeneous spaces}}{\placedmargin{extract a lemma on \unsure{morphisms} of \unsure{h}omogeneous spaces}}{marginal note placed}{ega25_25_180-v0002}
\var{\xout{$u\unsure{\in}Gl(\underline{E})(k)$}}{}{typewriter strike-out}{ega25_25_180-v0003}$u\hand{\in}Gl(\underline{E})(k)$, i.e. for every automorphism u of $\underline{E}$, u(U) will then
satisfy the same condition, and it suffices\hand{$\unsure{\leftarrow}$}\var{\inserted[typed]{to establish flatness}}{to establish flatness}{insertion}{ega25_25_180-v0004} to prove\hand{,} \var{\xout{that the u(}}{}{typewriter strike-out}{ega25_25_180-v0005}\hand{$\downarrow$}
% NOTE: The typed insertion has a long pen leader pointing left to the blotted mark after suffit and a downward branch after the canceled words; the curved ends are part of this leader.
of $(***)$\hand{,} that the u(U) cover $\underline{Q}(\underline{E})^*$ . Now for this it suffices to prove
that $Gl(\underline{E})(k)=$ Aut$(\underline{E})$ acts transitively on $\underline{Q}(\underline{E})^*(k)$, the set of
quadratic forms \var{\xout{on E}}{}{typewriter strike-out}{ega25_25_180-v0006} on $\overset{\hand{\checkmark}}{E}$ with nonzero discriminant, \var{\xout{\unsure{and}}}{}{typewriter strike-out}{ega25_25_180-v0007} i.e. that
\unsure{\var{\deleted{(}}{}{deletion}{ega25_25_180-v0008}}$***)$ is surjective, which will \unsure{th}\unsure{en} prove 16.11. We are thus reduced to
proving the

\uline{Lemma} 16.11.1. Let Q be a quadratic form in n variables over an
alg closed field k, such that $\overset{\hand{\sim}}{d}(Q)\hand{\neq} 0$ ; then Q is isomorphic to the standard form
% NOTE: Small irregular ink specks occur above the end of alors and the following Q in the source; no coherent accent or character can be established.
(i.e. can be put in standard form by a suitable choice of basis).

When E has even rank or k has char $\hand{\neq} 2$ , the hypothesis
means that Q is nondegenerate, and the conclusion is in Bourbaki
In the other case (k of char 2, \var{\xout{dim}}{}{typewriter strike-out}{ega25_25_180-v0010} rank E odd), one knows that Q is
degenerate, \var{\xout{therefore}}{}{typewriter strike-out}{ega25_25_180-v0011} let $E_1$ be a line in E lying in the kernel of the
associated form, $E_2$ a complement, so that Q takes the form

\mspage{\hand{180}}
\tnote{25-180-1}{The two parts of generic retain the a/i and q/i doubts in the source word; preceding retains its é/s doubt. The two parts of then retain the a/e and l/n doubts in alors. The retouched q in qui has no English spelling effect in which, but its old \unsure{q} remains doubtful. Alg closed means algebraically closed, and char means characteristic. The broken order around the typed flatness insertion, its two leader arrows, and the canceled opening parenthesis remain as recorded.}
